../cictikz/src/cictikz/data/tex/cictikz_preamble.tex

% A BPSK transmitter: the LO drives a differential pair, and the b0
% switches steer the current into the tank with either polarity.  The
% tank is coupled to the antenna through the matching inductor.  To the
% right, the symbols on the real and the imaginary axis.

\begin{circuitikz}[american, thick, transform shape, circuit ee IEC]

  ../cictikz/src/cictikz/data/tex/cictikz_lib.tex

  \def\xl{1.8}   % the left tank node
  \def\xr{6.2}   % the right tank node

  % the tank, and the matching inductor to the antenna
  \draw (\xl,5.4) to [L] (\xr,5.4);
  \draw (2.4,6.3) to [L] (5.6,6.3);
  \draw (2.4,6.3) -- (1.4,6.3) -- (1.4,6.1);
  \draw (1.4,6.1) \vground;
  \draw (5.6,6.3) -- (7.0,6.3) -- (7.0,7.1);
  \draw (6.6,7.7) -- (7.0,7.1) -- (7.4,7.7);

  % the four steering switches
  \draw (1.2,2.6) \lvnmos{MA}{};
  \draw (2.4,2.6) \lvmnmos{MB}{};
  \draw (5.6,2.6) \lvnmos{MC}{};
  \draw (6.8,2.6) \lvmnmos{MD}{};

  % the drains, tied to the tank nodes
  \draw (1.2,4.2) -- (1.2,4.6) -- (2.4,4.6) -- (2.4,4.2);
  \draw (\xl,4.6) -- (\xl,5.4);
  \fill (\xl,4.6) circle (0.075);
  \draw (5.6,4.2) -- (5.6,4.6) -- (6.8,4.6) -- (6.8,4.2);
  \draw (\xr,4.6) -- (\xr,5.4);
  \fill (\xr,4.6) circle (0.075);

  % the gates, b0 and its complement on both sides
  \draw (MA.gate) -- (0.5,3.4);
  \node[anchor=east] at (0.35,3.4) {\color{red}$b_0$};
  \draw (MB.gate) -- (3.1,3.4);
  \node[anchor=west] at (3.25,3.4) {\color{red}$\overline{b_0}$};
  \draw (MC.gate) -- (4.9,3.4);
  \node[anchor=east] at (4.75,3.4) {\color{red}$\overline{b_0}$};
  \draw (MD.gate) -- (7.5,3.4);
  \node[anchor=west] at (7.65,3.4) {\color{red}$b_0$};

  % the sources, cross coupled down to the two tail devices
  \draw (1.2,2.6) -- (1.2,2.0) -- (\xl,2.0);
  \draw (6.8,2.6) -- (6.8,2.0) -- (\xr,2.0);
  \draw (2.4,2.6) -- (\xr,2.0);
  \draw (5.6,2.6) -- (\xl,2.0);
  \fill (\xl,2.0) circle (0.075);
  \fill (\xr,2.0) circle (0.075);

  % the LO differential pair
  \draw (\xl,0.4) \lvnmos{ML}{};
  \draw (\xr,0.4) \lvmnmos{MR}{};
  \draw (\xl,0.4) -- (\xl,0.1);
  \draw (\xl,0.1) \vground;
  \draw (\xr,0.4) -- (\xr,0.1);
  \draw (\xr,0.1) \vground;

  % the LO, driving the two tails
  \draw (-1.9,0.4) rectangle (-0.4,2.0);
  \node[anchor=center] at (-1.15,1.2) {LO};
  \draw (-0.4,1.5) -- (\xl-0.5,1.5) -- (\xl-0.5,1.2);
  \draw (-0.4,0.9) -- (-0.1,0.9) -- (-0.1,-0.7) -- (\xr+0.5+0.8,-0.7)
        -- (\xr+0.5+0.8,1.2) -- (\xr+0.5,1.2);

  % the symbols on the real axis
  \def\px{9.4}
  \draw[-latex] (\px,4.8) -- (\px+4.0,4.8);
  \draw (\px,4.0) -- (\px,6.0);
  \node[anchor=east] at (\px-0.15,5.9) {$R$};
  \node[anchor=north] at (\px+3.95,4.7) {$t$};

  % the symbols on the imaginary axis
  \draw[-latex] (\px,2.2) -- (\px+4.0,2.2);
  \draw (\px,2.2) -- (\px,3.4);
  \node[anchor=east] at (\px-0.15,3.3) {$j$};
  \node[anchor=north] at (\px+3.95,2.1) {$t$};

  % the symbol instants
  \foreach \k/\yr in {1/5.6,2/4.0,3/5.6,4/4.0} {
    \draw[dotted] (\px+0.8*\k,5.8) -- (\px+0.8*\k,2.2);
    \fill[red] (\px+0.8*\k,\yr) circle (0.09);
    \fill[red] (\px+0.8*\k,2.2) circle (0.09);
  }
  \node[anchor=west] at (\px+4.4,5.7) {\color{red}1 \; 0 \; 1 \; 0};

\end{circuitikz}
\end{document}
